\textbf{Phần 1. Trắc nghiệm lựa chọn}
\setcounterpageref{ex}{0}

\begin{ex} %[2D8V39]
Kết quả của phép tính $2 \times 8$ là
\choice
{\True $16$}
{$14$}
{$18$}
{$15$}
\loigiai{
$2 \times 8 = 16$.
\textbf{Đáp án của câu hỏi là A}
}
\end{ex}

\begin{ex} %[2D8V39]
Một chuồng nuôi $7$ con gà. Hỏi đàn gà đó có tất cả bao nhiêu cái chân?
\choice
{$12$ cái chân}
{$16$ cái chân}
{\True $14$ cái chân}
{$10$ cái chân}
\loigiai{
Mỗi con gà có 2 cái chân: $2 \times 7 = 14$ (cái chân).
\textbf{Đáp án của câu hỏi là C}
}
\end{ex}

\begin{ex} %[2D8V39]
Tính giá trị biểu thức: $2 \times 9 - 8$.
\choice
{$9$}
{\True $10$}
{$11$}
{$8$}
\loigiai{
$2 \times 9 = 18$; $18 - 8 = 10$.
\textbf{Đáp án của câu hỏi là B}
}
\end{ex}

\begin{ex} %[2D8V39]
Mỗi bạn được phát $2$ quyển vở. Hỏi $6$ bạn được phát bao nhiêu quyển vở?
\choice
{$10$ quyển}
{$14$ quyển}
{$16$ quyển}
{\True $12$ quyển}
\loigiai{
$2 \times 6 = 12$ (quyển).
\textbf{Đáp án của câu hỏi là D}
}
\end{ex}

\begin{ex} %[2D8V39]
Tìm $x$ biết: $x : 2 = 9$.
\choice
{\True $18$}
{$16$}
{$17$}
{$20$}
\loigiai{
$x = 9 \times 2 = 18$.
\textbf{Đáp án của câu hỏi là A}
}
\end{ex}

\begin{ex} %[2D8V39]
Điền số thích hợp: $2 \times 5 + 2 \times 3 = 2 \times \dots$
\choice
{$7$}
{$9$}
{\True $8$}
{$6$}
\loigiai{
$2 \times (5 + 3) = 2 \times 8$. Điền 8.
\textbf{Đáp án của câu hỏi là C}
}
\end{ex}

\begin{ex} %[2D8V39]
Một đôi đũa gồm $2$ chiếc đũa. Hỏi $9$ đôi đũa có tất cả bao nhiêu chiếc đũa?
\choice
{$16$ chiếc}
{\True $18$ chiếc}
{$20$ chiếc}
{$14$ chiếc}
\loigiai{
$2 \times 9 = 18$ (chiếc).
\textbf{Đáp án của câu hỏi là B}
}
\end{ex}

\begin{ex} %[2D8V39]
Tính tổng: $2 + 2 + 2 + 2 + 2 + 2 + 10$.
\choice
{$20$}
{$24$}
{$18$}
{\True $22$}
\loigiai{
$2 \times 6 + 10 = 12 + 10 = 22$.
\textbf{Đáp án của câu hỏi là D}
}
\end{ex}

\begin{ex} %[2D8V39]
Lan có $5$ cái kẹo, Mai có số kẹo gấp đôi số kẹo của Lan. Hỏi cả hai bạn có tất cả bao nhiêu cái kẹo?
\choice
{\True $15$ cái kẹo}
{$10$ cái kẹo}
{$12$ cái kẹo}
{$14$ cái kẹo}
\loigiai{
Mai có: $5 \times 2 = 10$ (cái kẹo).\nCả hai bạn có: $5 + 10 = 15$ (cái kẹo).
\textbf{Đáp án của câu hỏi là A}
}
\end{ex}

\begin{ex} %[2D8V39]
Điền dấu thích hợp: $2 \times 6 \dots 2 \times 4 + 5$.
\choice
{$>$}
{$=$}
{\True $<$}
{Không so sánh được}
\loigiai{
Vế trái: $2 \times 6 = 12$. Vế phải: $8 + 5 = 13$. Vì $12 < 13$ nên điền dấu $<$. 
\textbf{Đáp án của câu hỏi là C}
}
\end{ex}

\begin{ex} %[2D8V39]
Có $8$ xe máy, mỗi xe máy có $2$ bánh xe. Hỏi $8$ xe máy có bao nhiêu bánh xe?
\choice
{$14$ bánh xe}
{\True $16$ bánh xe}
{$18$ bánh xe}
{$12$ bánh xe}
\loigiai{
$2 \times 8 = 16$ (bánh xe).
\textbf{Đáp án của câu hỏi là B}
}
\end{ex}

\begin{ex} %[2D8V39]
Tìm một số biết rằng lấy số đó trừ đi $6$ thì bằng $2 \times 7$.
\choice
{$18$}
{$19$}
{$21$}
{\True $20$}
\loigiai{
$2 \times 7 = 14$. Số đó: $14 + 6 = 20$.
\textbf{Đáp án của câu hỏi là D}
}
\end{ex}

\textbf{Phần 2. Trắc nghiệm đúng sai}
\setcounterpageref{ex}{0}

\begin{ex} %[2D8V39]
Xét tính đúng sai của các phép tính trong bảng nhân 2:
\choiceTF
{\True Mọi kết quả trong bảng nhân 2 đều là các số chẵn.}
{Phép tính 2 x 7 có kết quả lớn hơn 2 x 8.}
{\True Hai tích liền kề trong bảng nhân 2 luôn hơn kém nhau 2 đơn vị.}
{\True Biểu thức 2 x 4 + 2 bằng 2 x 5.}
\loigiai{
\begin{itemize}
    \item Ý a) Đúng: Số 2 nhân với bất kì số tự nhiên nào cũng được số chẵn. Đúng.
    \item Ý b) Sai: $2 \times 7 = 14 < 2 \times 8 = 16$.
    \item Ý c) Đúng: Đúng.
    \item Ý d) Đúng: $8 + 2 = 10 = 2 \times 5$. Đúng.
\end{itemize}
\begin{itemize}
    \item a) \textbf{Đúng}.
    \item b) \textbf{Sai}.
    \item c) \textbf{Đúng}.
    \item d) \textbf{Đúng}.
\end{itemize}
}
\end{ex}

\begin{ex} %[2D8V39]
Một đàn vịt có 8 con đang bơi dưới nước:
\choiceTF
{\True Đàn vịt có tất cả 16 cái chân.}
{\True Đàn vịt có tất cả 16 cái cánh.}
{Đàn vịt có tất cả 16 cái mỏ.}
{Nếu có thêm 2 con vịt nữa thì đàn vịt có 18 cái chân.}
\loigiai{
\begin{itemize}
    \item Ý a) Đúng: $2 \times 8 = 16$ cái chân.
    \item Ý b) Đúng: Mỗi con vịt có 2 cái cánh: $2 \times 8 = 16$ cái cánh.
    \item Ý c) Sai: Mỗi con vịt chỉ có 1 cái mỏ nên có 8 cái mỏ.
    \item Ý d) Sai: Thêm 2 con thì có 10 con: $2 \times 10 = 20$ cái chân (không phải 18).
\end{itemize}
\begin{itemize}
    \item a) \textbf{Đúng}.
    \item b) \textbf{Đúng}.
    \item c) \textbf{Sai}.
    \item d) \textbf{Sai}.
\end{itemize}
}
\end{ex}

\textbf{Phần 3. Trả lời ngắn}
\setcounterpageref{ex}{0}

\begin{ex} %[2D8V39]
Tính kết quả của phép nhân: 2 x 9.
\shortans{18}
\loigiai{
\textbf{Đáp số:} $18$\\
$2 \times 9 = 18$.
}
\end{ex}

\begin{ex} %[2D8V39]
Một con chim có 2 cái cánh. Hỏi đàn chim gồm 8 con có tất cả bao nhiêu cái cánh?
\shortans{16}
\loigiai{
\textbf{Đáp số:} $16$\\
$2 \times 8 = 16$ cái cánh.
}
\end{ex}

\begin{ex} %[2D8V39]
Tìm x biết: x - 2 x 6 = 8.
\shortans{20}
\loigiai{
\textbf{Đáp số:} $20$\\
$2 \times 6 = 12 \Rightarrow x = 8 + 12 = 20$.
}
\end{ex}

\begin{ex} %[2D8V39]
Một lớp học có 9 bàn đôi, mỗi bàn có 2 học sinh ngồi. Hỏi lớp học có bao nhiêu học sinh?
\shortans{18}
\loigiai{
\textbf{Đáp số:} $18$\\
$2 \times 9 = 18$ học sinh.
}
\end{ex}

\begin{ex} %[2D8V39]
Tính giá trị của biểu thức: 2 x 7 + 16.
\shortans{30}
\loigiai{
\textbf{Đáp số:} $30$\\
$2 \times 7 = 14$; $14 + 16 = 30$.
}
\end{ex}

\begin{ex} %[2D8V39]
Mai có 4 cái bút chì, số bút chì của Lan gấp đôi số bút chì của Mai. Hỏi cả hai bạn có bao nhiêu cái bút chì?
\shortans{12}
\loigiai{
\textbf{Đáp số:} $12$\\
Lan có: $4 \times 2 = 8$. Cả hai bạn: $4 + 8 = 12$ chiếc.
}
\end{ex}
